\subsection{The evaluated original baseline and its finite arithmetic receiver}

The canonical source still has its original order $k=4l+1$,
$e=1+k(m-1)$, $q=e(k+1)^2=2n$, and centre $c=k/2$.
Its multiplying polynomial is the original $f_l$, and its full
baseline remains $\mathcal B_k^0$. The following calculation uses
the complete BSL1--22 and HBL1--18 proofs included below.

Write $\gamma_j=\sqrt{2\pi}(2j)!/4^j$ for the original monic
Gamma norms and $\mu_q^{\rm even}=\int_{\mathbb R}y^{2q}d\sigma(y)$.
This last scalar is BSL's $\mu_q$ and the literal-index moment
$m_{2q}$ of LET. With the complete square pushforward
$\nu=(y\mapsto y^2)_*\sigma$, retain
\[
\begin{gathered}
 H_s^{(a)}=\det\left[\int_0^\infty t^{a+i+j}d\nu(t)\right]_{i,j<s},
 \qquad H_0^{(a)}=1,\\
 S_q=-\log\gamma_q-2\sum_{j=q+1}^{2q-1}\log\gamma_j
                                             -\log\gamma_{2q},\\
 Z_q=H_n^{(q)}H_{n+1}^{(q)}(H_n^{(q+1)})^2,
 \qquad V_q=S_q+\log Z_q-\log\mu_q^{\rm even}.
\end{gathered}\tag{BRI1}
\]
The original relation Grams are
$A_{\chi,r}=\operatorname{Gram}_\sigma[\chi S^j]_{j<r}$;
the power Grams are
$(M_{q,r})_{ij}=\int y^{2q+i+j}d\sigma(y)$.
Set $\eta_r=\log\det A_{\chi,r}-\log\det M_{q,r}$ for
$r=1,q,q+1$. The exact remainder map $J_N$ and its monic
relation frame give $\det G_N^\sigma=\det H_N^\sigma/
\det A_{\chi,N+1-q}$ at the four original ranks $0,1,q,q+1$.
Evenness of $\sigma$ splits the two high power Grams into precisely
the four factors in $Z_q$, while their scalar low Gram is
$\mu_q^{\rm even}$. Thus
\[
 \mathcal B_k^\sigma=V_q+\eta_q+\eta_{q+1}-\eta_1,
 \qquad \mathcal H_k=\mathcal B_k^\sigma-\mathcal B_k^0.
 \tag{BRI2}
\]
The second equality is the exact original source transition
$\mathcal B_k^0=\mathcal B_k^\sigma+\mathfrak T_k$ and
$\mathcal H_k=-\mathfrak T_k$. It identifies the scalar to
which the baseline bounds must be returned.

On the original full-source threshold
$k\geq2048\sqrt{d+g}$, retain the BSL constants
\[
\begin{gathered}
 a_k=\frac{k(k+2)(d-g)}{3q}<0,\quad
 t_\epsilon=\frac{k^2(d+g)}{q^2\epsilon^2}\leq\tfrac14,
 \quad d_\chi=\frac{qt_\epsilon^2}{2(1-t_\epsilon)},\\
 b_-=a_k\epsilon^{-2}-d_\chi,
 \qquad b_+=a_k64^{-2}+d_\chi,
 \quad L_*=L_q^{\rm tail}+L_{q+1}^{\rm tail}+L_1^{\rm tail},\\
 \mathsf L_k^\sigma=V_q+(2q+1)b_- -b_+-L_*,\quad
 \mathsf U_k^\sigma=V_q+(2q+1)b_+ -b_-+L_*,\\
 \mathsf L_k^\sigma\leq\mathcal B_k^\sigma
                                  \leq\mathsf U_k^\sigma.
\end{gathered}\tag{BRI3}
\]
Here $d=\delta^2$, $g=\gamma^2$, $\epsilon=2^{-10}$ and
all three $L_r^{\rm tail}$ are the original complete tail constants
already displayed above; $b_-,b_+$ are the packet density bounds,
while $b_k^-,b_k^+$ keep their different original high-error meaning.
The exact multiplication map between the bulk sources is
$P(c+iy)(iy)^q\mapsto P(c+iy)\chi(c+iy)$.
Its squared density ratio is between $\exp(b_-)$ and
$\exp(b_+)$, by the full paired-root logarithm in BSL6--8.
Congruence on the same relation coefficient space and both positive
full-source tails gives $rb_- -L_r^{\rm tail}\leq\eta_r
\leq rb_++L_r^{\rm tail}$. Taking the two high positive signs
and the scalar low negative sign proves BRI3 from BRI2.

By the second equality in BRI2 the interval
$[\mathsf L_k^\sigma-\mathcal B_k^0,
\mathsf U_k^\sigma-\mathcal B_k^0]$ contains the same original
$\mathcal H_k$ as every previously retained Gamma interval.
It can therefore be included in their exact maximum/minimum
intersection. No change is made to any original source or
arithmetic allowance. Independently of which finite interval is
tightest, the original equality
$\mathcal B_k^{\rm ar}=\mathcal B_k^\sigma+\delta_k^\sigma$
gives the explicit finite baseline receiver
\[
\begin{gathered}
 \mathsf L_k^\sigma-a_k^-\leq\mathcal B_k^{\rm ar}
                                   \leq\mathsf U_k^\sigma+a_k^+,\\
 \mathsf L_k^\sigma-D_k^-\leq\mathcal B_k^{\rm ar}
                                   \leq\mathsf U_k^\sigma+D_k^+,\\
 \max\{0,\mathsf L_k^\sigma-D_k^- -4q\log(D_hk)\}
 \leq\mathcal R_k
 \leq\mathsf U_k^\sigma+D_k^+ -4q\log(D_hk).
\end{gathered}\tag{BRI4}
\]
The first line adds the unchanged
$\delta_k^\sigma\in[-a_k^-,a_k^+]$ to BRI3; the second
uses $a_k^\pm\leq D_k^\pm$. The last subtracts the unchanged
HC threshold and intersects the lower endpoint with the original
$\mathcal R_k\geq0$, on the combined original source and HC domains.

The complete HBL13--18 coefficient is
\[
 C_B=9-8\log2+\mathcal F(2,\pi),\qquad
 C_B>\frac{769}{17010}>0,\qquad
 \frac{V_q}{q^2}\longrightarrow C_B.
 \tag{BRI5}
\]
Here $\mathcal F(2,\pi)$ is the supremum of the single-kernel
energy for $\pi\sqrt x-2\log x$, attained by the exact positive
density $\rho_{2,\pi}$ already specified in GRI2. HBL14--16
gives its endpoint-integral expression, and HBL17--18 proves
the strict rational lower bound by its explicit arcsine competitor.
The determinant calculation retains the complete factors
\[
\begin{split}
 S_q&=-6q^2\log q+(9-8\log2)q^2+O(q\log q),\\
 \log Z_q&=(6q^2+2q)\log q+q^2\mathcal F(2,\pi)+o(q^2),\qquad
 \log\mu_q^{\rm even}=O(q\log q).
\end{split}
\]
Their $q^2\log q$ terms cancel only in BRI1. BSL6--10 gives
$\eta_q+\eta_{q+1}-\eta_1=O(q/e)=o(q^2)$, while BRI3
retains its finite asymmetric enclosure.

Consequently the full original baselines, arithmetic volume and
HC residual have the evaluated common leading value
\[
\begin{gathered}
 \frac{\mathcal B_k^\sigma}{q^2}\longrightarrow C_B,\qquad
 \frac{\mathcal B_k^0}{q^2}\longrightarrow C_B,\qquad
 \frac{\mathcal B_k^{\rm ar}}{q^2}\longrightarrow C_B,\qquad
 \frac{\mathcal R_k}{q^2}\longrightarrow C_B,\\
 \frac{\mathcal B_k^0-\mathcal B_k^{\rm ar}}{kq}
                \longrightarrow\frac{\mathcal C_\Gamma}{4},
 \qquad \frac{\mathcal B_k^0}{kq}\longrightarrow+\infty.
\end{gathered}\tag{BRI6}
\]
To prove these statements, BRI2--5 first gives the fixed Gamma
baseline. The original tensor return is
$\mathcal B_k^0=\mathcal B_k^\sigma+\mathfrak T_k$ with
$\mathfrak T_k/(lq)\to\mathcal C_\Gamma$ and $l/q\to0$.
The arithmetic return is
$\mathcal B_k^{\rm ar}=\mathcal B_k^\sigma+\delta_k^\sigma$,
where the retained simple and multiple branches prove
$\delta_k^\sigma/(kq)\to0$ and $k/q\to0$.
The original threshold satisfies $4\log(D_hk)/q\to0$.
These give the first line. The finer difference in the second line
is GRI5--7, unchanged. Finally $q/k\to\infty$ and $C_B>0$
prove the displayed divergence. In particular the previous necessary
lower bound at the $kq$ scale remains valid and is now strengthened
by the evaluated $q^2$ baseline; no finer remainder is assumed.

The original HC identity retains its actual components:
\[
 \mathcal B_k^{\rm ar}=4q\log\mathcal L_k-W_k^\sigma
                  +\mathcal S_k-\tau_k^{\rm win},\qquad
 \mathcal L_k=2\delta e(k+1)\lfloor(k+1)^2/4\rfloor.
 \tag{BRI7}
\]
The symbol $W_k^\sigma$ here is the original factorial term of
TWA5, and is distinct from the universal Gamma centre $W_k$.
Its complete formula gives $O(q\log q)$; TWA6 gives
$\tau_k^{\rm win}=O_h(k+\log q)$. Also
$\log\mathcal L_k=O_h(\log k)$ and $\log q=O_m(\log k)$.
Dividing the exact identity by $q^2$ and using BRI6 proves
\[
 \frac{\mathcal S_k}{q^2}\longrightarrow C_B>0.
 \tag{BRI8}
\]
Thus the original contraction, phase and relative-control sum has
an evaluated cost on this scale, with its individual terms retained.

Finally let $K_{\rm ker}=\ker\Lambda$ and retain the original
coefficient section $S_*:B_{\rm obs}\to E_\chi$.
For each of the same positive quotient metrics, the full frame map
$[I_{K_{\rm ker}},S_*]:K_{\rm ker}\oplus B_{\rm obs}\to E_\chi$
is an isomorphism, and the exact block determinant gives
\[
 \det(I_{K_{\rm ker}}^*G_NI_{K_{\rm ker}})\det Q_N^{B_{\rm obs}}
 =|\det[I_{K_{\rm ker}},S_*]|^2\det G_N.
 \tag{BRI9}
\]
The inverse splitting is $z\mapsto(z-S_*\Lambda z,\Lambda z)$;
orthogonal block elimination proves the determinant identity.
The frame factor cancels with signs $+,+,-,-$ at the same four
endpoints. BSL20--21 therefore transports the same result to the
actual combined mixed observation:
$(F_K(1)+F_B(1))/q^2\to C_B$. Each individual kernel and
boundary Gram, its off-diagonal entries and its Taylor unit
remain present. No individual factor is assigned the combined limit.


\subsection{The refined original low endpoint in the same complete return}

Keep the original $k=4l+1$, $q=e(4l+2)^2$ and both original
high ranks. The complete GDR1--22 and LER1--23 proofs give,
with all moment indices retained,
\[
 \log\frac{m_{2q+2l}}{m_{2q}}
 =F_{q,l}-K_{q,l}+r_{q,l}+\varepsilon_{q,l},\qquad
 -T_{q,l}\leq r_{q,l}\leq0,\quad
 |\varepsilon_{q,l}|\leq E_{q,l},
 \tag{LRI1}
\]
where the complete finite quantities are
\[
\begin{gathered}
 F_{q,l}=2l\log\frac{4q}{\pi}+\frac{l^2}{q}
 -\frac{l(16l^2-1)}{48q^2}
 +\frac{l^2(8l^2-1)}{48q^3},\\
 K_{q,l}=\frac{\pi^2}{64}\left\{
 \frac1{(2q-3/2)(2q-1/2)}
 -\frac1{(2q+2l-3/2)(2q+2l-1/2)}\right\}>0,\\
 T_{q,l}=\frac{l(768l^4-160l^2+7)}{7680q^4},\qquad
 E_{q,l}=\frac{51l}{16q^5}.
\end{gathered}\tag{LRI2}
\]
The literal-index ratio here equals LER's even-index ratio
$\mu_{q+l}/\mu_q$; the whole mass $\sqrt{2\pi}$ and both
half-axes remain present before cancellation. GDR proves the
global density identity with its explicit remainder, and LER3--14
integrates and telescopes it at every original integer exponent.
The exact cubic logarithm identity then gives the one-sided
$r_{q,l}$ in LER18. Thus no formal asymptotic expansion is
substituted for the finite equality LRI1.

The paired high determinants remain exactly
$\rho_q+\rho_{q+1}=\log(Z_{q+l}/Z_q)$, with the original
$Z_a=H_n^{(a)}H_{n+1}^{(a)}(H_n^{(a+1)})^2$ and $n=q/2$.
Define the refined, still complete high-determinant expression
\[
 \widehat W_k=P_k+\rho_q+\rho_{q+1}-F_{q,l}+K_{q,l}.
\]
Substitution in the original definition of $W_k$ gives
\[
 W_k=\widehat W_k-r_{q,l}-\varepsilon_{q,l},\qquad
 \widehat W_k-E_{q,l}\leq W_k
                  \leq\widehat W_k+T_{q,l}+E_{q,l}.
 \tag{LRI3}
\]
In particular the rational Gamma correction enters $W_k$
positively and the $l^2/q$ term enters negatively. Both signs
come from the original low-endpoint subtraction.

Use the same source errors $e_0\in[-a_0,b_0]$ and
$e_q+e_{q+1}\in[-b_k^+,b_k^-]$ in
$\mathcal H_k=-W_k+e_0-e_q-e_{q+1}$. This proves
\[
\begin{gathered}
 -\widehat W_k-T_{q,l}-E_{q,l}-b_k^--a_0
 \leq\mathcal H_k
 \leq-\widehat W_k+E_{q,l}+b_k^++b_0,\\
 \mathcal B_k^0-\widehat W_k-T_{q,l}-E_{q,l}-b_k^--a_0-a_k^-
 \leq\mathcal B_k^{\rm ar}\\
 \leq\mathcal B_k^0-\widehat W_k+E_{q,l}+b_k^++b_0+a_k^+.
\end{gathered}\tag{LRI4}
\]
For the second enclosure add the unchanged arithmetic interval;
its established outer version replaces $a_k^\pm$ by $D_k^\pm$.
Subtracting the original HC threshold and imposing the original
nonnegative lower constraint gives its corresponding HC enclosure.
For fixed exact $W_k$, LRI3 proves that LRI4 contains the earlier
WRC5 interval, so it does not remove the retained tighter finite
intersection. Its new content is the exact original low correction
and its finite signed error inside the complete high-determinant
expression. No high determinant or source unit has been removed.

Finally LER23 proves
$\log(m_{2q+2l}/m_{2q})-2l\log(4q/\pi)-l^2/q\to0$.
The term $l^2/q$ tends to $1/16$ at $m=1$ and to zero for
each fixed $m\geq2$. These are two original multiplicity
branches of this scalar endpoint, retained alongside the full
$W_k$ and both of its high relation ranks.


\subsection{The finite four-block estimate returned to the original signed maps}

The complete QLG1--32 and QGT1--28 proofs retain the original
$k=4l+1$, $q=[1+k(m-1)](k+1)^2=2n$ and the original Gamma
density. QGT15 is the exact sum $\mathcal C_{q,l}$ of the four
free-energy increments at ranks $n,n+1,n,n$ with multiplicities
$1,1,1,1$ and their actual quarter-shifted parameters. With its
proved finite constants $A=M/2+2C$, define
\[
\begin{gathered}
 \mathcal U_{q,l}=Al(3\sqrt n+\sqrt{n+1}),\\
 \mathcal V_{q,l}=\widetilde P_{q,l}+\mathcal C_{q,l}
                         +2l\log q-F_{q,l}+K_{q,l},\\
 \omega_{q,l}^-=\mathcal V_{q,l}-\mathcal U_{q,l}-E_{q,l},
 \qquad
 \omega_{q,l}^+=\mathcal V_{q,l}
                        +\mathcal U_{q,l}+T_{q,l}+E_{q,l}.
\end{gathered}\tag{QRI1}
\]
Here $\widetilde P_{q,l}=P_k+4lq\log q$ is the exact original
finite source sum, while $F,K,T,E$ are precisely LRI2. The
quantile-cell and circle-potential proof QLG, transported through
the exact square pushforward QGT3--11, gives a finite $Cs$ error
at each determinant rank $s$. The convex secant proof QGT12--17
then gives the exact full high ratio
\[
 \log(Z_{q+l}/Z_q)=(4q+2)l\log q+\mathcal C_{q,l}+\xi_{q,l},
 \qquad |\xi_{q,l}|\le\mathcal U_{q,l}.
\]
Substituting LRI1, with its original negative low-endpoint sign,
gives
$W_k=\mathcal V_{q,l}+\xi_{q,l}-r_{q,l}-\varepsilon_{q,l}$.
Thus $\omega_{q,l}^-\le W_k\le\omega_{q,l}^+$.

There is a second finite enclosure with no finite source sum in
its centre. Keep the explicit $p^{\rm c}_{q,l},p^-_{q,l},p^+_{q,l}$
and $\mathcal V^{\rm c}_{q,l}$ of QGT21, including all their
linear, quadratic and cubic product corrections. Set
\[
\begin{gathered}
 \omega_{q,l}^{\rm c,-}=\mathcal V^{\rm c}_{q,l}
                    -p^-_{q,l}-\mathcal U_{q,l}-E_{q,l},\\
 \omega_{q,l}^{\rm c,+}=\mathcal V^{\rm c}_{q,l}
                    +p^+_{q,l}+\mathcal U_{q,l}+T_{q,l}+E_{q,l}.
\end{gathered}\tag{QRI2}
\]
The proved original product interval
$-p^-_{q,l}\le\widetilde P_{q,l}-p^{\rm c}_{q,l}\le p^+_{q,l}$
implies
$\omega_{q,l}^{\rm c,-}\le\omega_{q,l}^-\le W_k
 \le\omega_{q,l}^+\le\omega_{q,l}^{\rm c,+}$.
This also proves the finite containment relation between these
two evaluated presentations.

For either pair $(\omega^-,\omega^+)$ in QRI1 or QRI2 the actual
source identity, with all original errors, is
\[
\begin{gathered}
 \mathcal H_k=-\mathfrak T_k=-W_k+e_0-e_q-e_{q+1},\\
 e_0\in[-a_0,b_0],\qquad
 e_q+e_{q+1}\in[-b_k^+,b_k^-],\\
 -\omega^+-b_k^--a_0\le\mathcal H_k
                            \le-\omega^-+b_k^++b_0.
\end{gathered}\tag{QRI3}
\]
The lower bound uses $-W_k\ge-\omega^+$, $e_0\ge-a_0$ and
$-(e_q+e_{q+1})\ge-b_k^-$; the upper bound uses the other three
endpoints. Thus every sign is fixed by the original map. This
interval contains the exact-$W_k$ interval WRC5. Its intersection
with the previously retained finite intervals therefore preserves
every previously proved tighter bound.

Adding the unchanged arithmetic scalar and then using its original
majorants gives, for either of the same pairs,
\[
\begin{split}
 \mathcal B_k^0-\omega^+-b_k^--a_0-a_k^-
 &\le\mathcal B_k^{\rm ar}
 \le\mathcal B_k^0-\omega^-+b_k^++b_0+a_k^+,\\
 \mathcal B_k^0-\omega^+-b_k^--a_0-D_k^-
 &\le\mathcal B_k^{\rm ar}
 \le\mathcal B_k^0-\omega^-+b_k^++b_0+D_k^+.
\end{split}\tag{QRI4}
\]
Indeed $\mathcal B_k^{\rm ar}=\mathcal B_k^0+\mathcal H_k
+\delta_k^\sigma$, $\delta_k^\sigma\in[-a_k^-,a_k^+]$ and
$0\le a_k^\pm\le D_k^\pm$. The original source and arithmetic
domains are intersected; no original degree or multiplicity
restriction is weakened. On the original HC domain the same map
$\mathcal R_k=\mathcal B_k^{\rm ar}-4q\log(D_hk)$ and the
proved $\mathcal R_k\ge0$ give
\[
\begin{split}
 \max\{0,\mathcal B_k^0-\omega^+-b_k^--a_0-D_k^-
                              -4q\log(D_hk)\}
 &\le\mathcal R_k\\
 &\le\mathcal B_k^0-\omega^-+b_k^++b_0+D_k^+
                              -4q\log(D_hk).
\end{split}\tag{QRI5}
\]

For the quantitative return around the previously evaluated coefficient,
let $\mathcal D_{q,l}$ be the full nonnegative finite expression QGT25
and write $\mathcal E_{q,l}^{\rm high}=\mathcal U_{q,l}+\mathcal D_{q,l}$.
QGT26 and the same source errors prove
\[
\begin{gathered}
 |W_k-lq\mathcal C_\Gamma|\le\mathcal E_{q,l}^{\rm high},\\
 -\mathcal E_{q,l}^{\rm high}-b_k^--a_0
 \le \mathcal H_k+lq\mathcal C_\Gamma
 \le \mathcal E_{q,l}^{\rm high}+b_k^++b_0.
\end{gathered}\tag{QRI6}
\]
The finite expression QGT25 and the exact packet relation give
$\mathcal E_{q,l}^{\rm high}=O(l\sqrt q+l^2)$.
The already proved original source errors are $O(k)$ for $m=1$
and $O(1)$ for each fixed $m\ge2$, and hence are $o(q)$ in both
branches. For each fixed $m\ge2$, $l/\sqrt q\to0$ and
$l^2/q\to0$; for $m=1$, $q=(4l+2)^2$ and both ratios remain
bounded. Dividing QRI6 by $q$ therefore proves
\[
 \begin{cases}
 (\mathcal H_k+lq\mathcal C_\Gamma)/q\longrightarrow0,
                   &m\ge2\text{ fixed},\\
 \mathcal H_k+lq\mathcal C_\Gamma=O(q),&m=1.
 \end{cases}
 \tag{QRI7}
\]
The first statement is equivalently
$(\mathfrak T_k-lq\mathcal C_\Gamma)/q\to0$ for the same fixed
multiple branch. The second keeps the entire finite uncertainty
in QRI3 and QRI6 and assigns no unproved limit to that quotient.

For the original arithmetic return the exact counterpart is
\[
\begin{gathered}
 \Delta_k^\Gamma+lq\mathcal C_\Gamma-\delta_k^\sigma
       =\mathcal H_k+lq\mathcal C_\Gamma,\\
 \mathcal B_k^{\rm ar}-\mathcal B_k^0+lq\mathcal C_\Gamma
       -\delta_k^\sigma=\mathcal H_k+lq\mathcal C_\Gamma,\\
 \mathcal R_k-\mathcal B_k^0+4q\log(D_hk)
       +lq\mathcal C_\Gamma-\delta_k^\sigma
       =\mathcal H_k+lq\mathcal C_\Gamma.
\end{gathered}\tag{QRI8}
\]
Each equality is the original definition followed by addition of
the displayed identical quantities. Each right-hand side obeys
QRI6 and both branches of QRI7. The arithmetic scalar remains
on the left with its full asymmetric bounds; in the multiple
branch the earlier $q\log k$ constants $64m+245$ and $128m+490$
remain unchanged. At $m=1$ its established $5/7$ rate also
remains unchanged. These equalities return the newly proved
$q$-scale Gamma information without assigning that scale to an
unestimated arithmetic contribution.

The original simple branch now has a more precise centre as well.
QGQ4--9 proves, from the complete QGT21 finite expression and its
actual four shifted equilibrium integrals,
\[
 \mathcal C_\Gamma^{[q,1]}
       =\frac{\partial_\alpha\mathcal L(2,\pi)}8
                            -\frac{\log2}{4},\qquad
 \frac{\mathcal V^{\rm c}_{q,l}-lq\mathcal C_\Gamma}{q}
                            \longrightarrow\mathcal C_\Gamma^{[q,1]}
 \quad(m=1).
\]
Its proof retains the Taylor remainder of every high block and
every original source and scalar correction. QGQ9 bounds the
remaining actual-$W_k$ error about that centre by $A/\sqrt2$
in the limit superior. Since
$(e_0-e_q-e_{q+1})/q\to0$ on the same original branch, the
identity in QRI3 gives the correctly signed return
\[
 \boxed{\displaystyle
 \limsup_{l\to\infty}
 \left|\frac{\mathcal H_k+lq\mathcal C_\Gamma}{q}
                         +\mathcal C_\Gamma^{[q,1]}\right|
                \le\frac A{\sqrt2}\qquad(m=1).}
 \tag{QRI9}
\]
Indeed the expression inside the absolute value equals the
negative of
$(W_k-lq\mathcal C_\Gamma)/q-\mathcal C_\Gamma^{[q,1]}$
plus the displayed vanishing source error. The triangle inequality
and QGQ9 prove QRI9. By each exact equality in QRI8 the same bound
applies to its left-hand side divided by $q$ and then increased
by $\mathcal C_\Gamma^{[q,1]}$. In particular the arithmetic scalar
is still subtracted there with its original sign. The finite
interval QRI3 remains valid before passing to this limit superior.

\subsection{The original Gamma source mapped into each mixed row}

The accepted GMB1--16 proof computes the original mixed rows from
the original quotient relation. Keep $s\in\{1,k\}$, the packet
centre $c=k/2$, its polynomial $\chi$, the fixed remainder frame
of $E_\chi=\mathbb C[S]/(\chi)$, and the same surjection
$\Lambda:E_\chi\to B_{\rm obs}$. The source order $s=1$ is
the original fixed Gamma comparison, and $s=k$ is the original
tensor Gamma source. The source norm is
$h_d^{(s)}=(2\pi)^{s/2}d!(s/2)_d$.

Let $T_\chi$ be multiplication by $S$ in the fixed remainder
frame, and put $Y_\chi=(T_\chi-cI)/i$. The finite original-source
recurrence is
\[
\begin{gathered}
 r_0=[1],\quad r_1=Y_\chi[1],\qquad
 r_{d+1}=Y_\chi r_d-d(s/2+d-1)r_{d-1},\\
 C_s=\left[\frac{r_0}{\sqrt{h_0^{(s)}}},\ldots,
          \frac{r_{q-1}}{\sqrt{h_{q-1}^{(s)}}}\right],\qquad
 v_j=\frac{C_s^{-1}r_{q+j}}{\sqrt{h_{q+j}^{(s)}}},\\
 K_j=I_q+\sum_{a<j}v_av_a^*,\qquad
 t_j=v_j^*K_j^{-1}v_j,\qquad R_s=\Lambda C_s,\\
 M_j^{\rm obs}=R_sK_jR_s^*,\qquad
 w_j=R_sv_j,\qquad
 \xi_j=w_j^*(M_j^{\rm obs})^{-1}w_j.
\end{gathered}\tag{MRI1}
\]
The conjugate transpose is denoted by a star. The recurrence is
the remainder of the original monic Gamma orthogonal polynomial
recurrence; its determinant is
$\det C_s=i^{-q(q-1)/2}\prod_{d<q}(h_d^{(s)})^{-1/2}$.
It therefore gives an invertible map with its actual source
phase and mass. Positivity of $K_j$ and surjectivity of $R_s$
prove the stated inverses exist. For zero boundary dimension
the matrices are empty and $\xi_j=0$.

The original remainder map at degree $q+j-1$, in these orthonormal
source coordinates, is precisely
$C_s[I_q,v_0,\ldots,v_{j-1}]$. Its minimum norm lift of $C_sx$
is $[I_q,v_0,\ldots,v_{j-1}]^*K_j^{-1}x$. The orthogonality to
the full relation kernel proves the actual quotient metric
$G_{q+j-1}^{(s)}=C_s^{-*}K_j^{-1}C_s^{-1}$.
Furthermore
\[
\begin{gathered}
 r_j^{\rm ker}=v_j-K_jR_s^*(M_j^{\rm obs})^{-1}w_j,
 \qquad R_sr_j^{\rm ker}=0,\\
 t_j=\xi_j+(r_j^{\rm ker})^*K_j^{-1}r_j^{\rm ker},
 \qquad0\le\xi_j\le t_j.
\end{gathered}\tag{MRI2}
\]
The two summands of $v_j$ are orthogonal in the $K_j^{-1}$
metric since their cross inner product contains $R_sr_j^{\rm ker}=0$.
The map $C_s$ restricts to a bijection
$\ker R_s\to\ker\Lambda$. Thus the original kernel and its
observation are both retained in the same computation.

Put $D=q+j$. The original monic relation has the exact squared
norm $a_{D-1}=h_D^{(s)}(1+t_j)$ and coefficient row
$\alpha_{D-1}=-i^{-D}\sqrt{h_D^{(s)}}\,
v_j^*K_j^{-1}C_s^{-1}$. The full section and block inverse
calculation in GMB11--13 gives the original two contractions
\[
\begin{gathered}
 \beta_{D-1}=h_D^{(s)}(1+t_j)\frac{t_j-\xi_j}{1+\xi_j},
 \qquad
 \nu_{D-1}=h_D^{(s)}(1+t_j)^2\frac{\xi_j}{1+\xi_j},\\
 1+\frac{\beta_{D-1}}{a_{D-1}}
          =\frac{1+t_j}{1+\xi_j},\qquad
 1+\frac{\nu_{D-1}}{a_{D-1}+\beta_{D-1}}=1+\xi_j.
\end{gathered}\tag{MRI3}
\]
These are the original BRU factors in their original order. Their
product is $1+t_j$, equal to
$\det G_{D-1}^{(s)}/\det G_D^{(s)}$ by the exact rank-one
inverse update. The terms in each denominator and the phase in
the row remain part of the original map.

Define the actual four-endpoint weight
$\mathcal W_q(f)=f_0+2\sum_{j=1}^{q-1}f_j+f_q$. Telescoping
the preceding quotient and mixed determinants gives
\[
\begin{gathered}
 \mathcal B_k^{(s)}=\mathcal W_q\bigl(\log(1+t_j)\bigr),\\
 F_K^{(s)}=\mathcal W_q\left(\log\frac{1+t_j}{1+\xi_j}\right),
 \quad F_B^{(s)}=\mathcal W_q\bigl(\log(1+\xi_j)\bigr),\\
 F_K^{(s)}+F_B^{(s)}=\mathcal B_k^{(s)},\qquad
 0\le F_K^{(s)}\le\mathcal B_k^{(s)},\quad
 0\le F_B^{(s)}\le\mathcal B_k^{(s)}.
\end{gathered}\tag{MRI4}
\]
The weights are exactly those of degrees $q-1,q,2q-1,2q$;
the relation at step $j$ has degree $q+j$. Positivity in MRI2
proves each displayed inequality term by term. BRI6 and BRI9
therefore apply to the complete sum in MRI4 for each of the
two original source orders. The finite formulas MRI1--3 give
the individual factors without assigning either one the combined
coefficient $C_B$.

The arithmetic transition remains the exact original identity
\[
 \mathcal B_k^{\rm ar}
   =F_K^{(k)}+F_B^{(k)}+\mathcal H_k+\delta_k^\sigma
   =F_K^{(1)}+F_B^{(1)}+\delta_k^\sigma.
 \tag{MRI5}
\]
Indeed MRI4 identifies the two sums with $\mathcal B_k^0$ and
$\mathcal B_k^\sigma$, respectively, while BRI2 gives
$\mathcal H_k=\mathcal B_k^\sigma-\mathcal B_k^0$.
This proves the exact map of the new row formulas to the same
arithmetic scalar and therefore to QRI4--8. Every primary
coefficient and Taylor unit in $\Lambda$ remains in $R_s$.
No new source order or multiplying factor is selected.

\subsection{The exterior Gamma law returned through its full relation mass}

The original input source remains $dm_{0,s}$ with $s\in\{1,k\}$,
mass $M_s=(2\pi)^{s/2}$, and parameter $b=s/2$. For the actual
relation degree $r\in\{1,q,q+1\}$ retain the exterior measure
\[
 d\lambda_{\chi,s,r}(y)=\frac1{r!}
       \prod_{i<j}(y_j-y_i)^2
       \prod_{i=1}^r|\chi(c+iy_i)|^2\,dm_{0,s}(y_i),
 \qquad \nu_{\chi,s,r}=(\Sigma_r)_*\lambda_{\chi,s,r},
 \quad \Sigma_r(y)=\sum_i y_i.
 \tag{IRI1}
\]
IGO7 proves that this is the squared modulus of the original
relation exterior vector, including its phase
$i^{r(r-1)/2}$ and factor $1/\sqrt{r!}$. Its total mass is the
original relation Gram determinant $Z_{s,r}$ of IFB17. At $r=0$
the space is the unit point, the measure is $\delta_0$, and the
mass is one.

For $r\ge1$ put $B=br+r(r-1)$. IGO9--11 constructs the exact
polynomial $D_{\chi,s,r}(z)=\sum_{j=0}^{qr}d_jz^{2j}$ by the
displayed differential recurrence from every coefficient of
$|\chi(c+iy)|^2$. Its full signed constituent and positive
polynomial formulas are
\[
\begin{gathered}
 a_h=\sum_{j=h}^{qr}(-1)^{j+h}\binom jh d_j,\qquad
 R_{\chi,s,r}(X)=\sum_{h=0}^{qr}a_h
       \frac{\Gamma(B)}{\Gamma(B+2h)}
       \prod_{j=0}^{h-1}[X+(B+2j)^2],\\
 d\nu_{\chi,s,r}(Y)
 =\sum_{h=0}^{qr}\frac{M_s^ra_h}{(2\pi)^{B+2h}}
                             dm_{0,2B+4h}(Y)
 =\frac{M_s^r}{(2\pi)^B}
                  R_{\chi,s,r}(Y^2)dm_{0,2B}(Y).
\end{gathered}\tag{IRI2}
\]
IGO15--21 proves the first equality by the finite binomial expansion
of $D(i\tanh\xi)$ in the exact characteristic transform; the Gamma
recurrence proves the second equality, with all powers of two
accounted for. The polynomial has degree $qr$ and is strictly
positive on $X\ge0$ by the original positive fibre integral.
The input source orders have not been selected anew. The produced
orders are $sr+2r(r-1)+4h$, with every $a_h$ retained. Each
constituent on its natural line $U_h=B+2h+iY$ is pulled back to
the actual sum line $T=rc+iY$ by
$U_h=T+(B+2h-rc)$. This translation and its inverse preserve every
polynomial degree and have coefficient determinant one.

The total mass, including all signed cancellations, is
\[
\begin{gathered}
 Z_{s,r}=M_s^rD_{\chi,s,r}(0)
        =M_s^r\sum_{h=0}^{qr}a_h,\\
 D_{\chi,s,r}(0)
   =\left(\prod_{j=0}^{r-1}(q+j)!(b)_{q+j}\right)\Omega_{s,r},
 \qquad
 \Omega_{s,r}=\det\left(I_q+\sum_{j<r}v_jv_j^*\right).
\end{gathered}\tag{IRI3}
\]
Here $v_j$ are precisely the original remainder coordinates MRI1.
To prove the second line, IFB9 identifies the full relation space
in orthonormal source coordinates as the graph $(-V_rz,z)$.
Its monic high-coordinate matrix is triangular with squared
determinant $\prod_{j<r}h_{q+j}^{(s)}$; the graph Gram is
$I_r+V_r^*V_r$. Taking determinants, using
$\det(I_r+V_r^*V_r)=\det(I_q+V_rV_r^*)$, and then substituting
$h_{q+j}^{(s)}=M_s(q+j)!(b)_{q+j}$ proves IRI3 from the first
line. Every source mass is present before this equality cancels it.
At $r=0$ both products and both determinants are one.

Return these actual masses to the same four original quotient
degrees. IFB30 and IRI3 give
\[
\begin{gathered}
 \mathcal B_k^{(s)}
   =S_q^{(s)}+\log Z_{s,q}+\log Z_{s,q+1}-\log Z_{s,1},\\
 S_q^{(s)}=-\log h_q^{(s)}
       -2\sum_{j=q+1}^{2q-1}\log h_j^{(s)}-\log h_{2q}^{(s)},\\
 \mathcal B_k^{(s)}
   =\log\Omega_{s,q}+\log\Omega_{s,q+1}-\log\Omega_{s,1}
   =\mathcal W_q\bigl(\log(1+t_j)\bigr).
\end{gathered}\tag{IRI4}
\]
The two high monic norm products divided by the scalar low norm
contain $h_q$ once, $h_{q+1},\ldots,h_{2q-1}$ twice, and
$h_{2q}$ once. These are exactly the factors with negative signs
in $S_q^{(s)}$; their full source mass powers are $2q$ and $-2q$.
This proves the last line before any asymptotic passage. The final
equality telescopes $\Omega_{s,j+1}/\Omega_{s,j}=1+t_j$ with
the same weight-one and weight-two endpoints as MRI4. Thus the
intrinsic exterior law returns to the actual baseline already
evaluated in BRI5--6, through the full original polynomial mass.
MRI3--5 then gives its individual mixed factors and arithmetic map.

For completeness the operation also retains its Hilbert-space
kernel. The pullback $U_rF=F\circ\Sigma_r$ is an isometry from
$L^2(\nu_{\chi,s,r})$ to $L^2(\lambda_{\chi,s,r})$. If
$w(y)$ is the complete Lebesgue density of IRI1 and
$p(Y)=\int_{\mathbb R^{r-1}}w(u,Y-\sum u)du$, its exact adjoint is
\[
 (U_r^*f)(Y)=\frac{\int_{\mathbb R^{r-1}}
                 f(u,Y-\sum u)w(u,Y-\sum u)\,du}{p(Y)},
 \qquad
 L^2(\lambda_{\chi,s,r})
       =U_rL^2(\nu_{\chi,s,r})\oplus\ker U_r^*.
 \tag{IRI5}
\]
IGO22--26 proves $p(Y)>0$ and the fibre Cauchy--Schwarz bound,
which proves the adjoint and the orthogonal decomposition. The
kernel is exactly the functions with zero displayed weighted fibre
integral almost everywhere. Multiplication by the original exterior
vector gives the isometry to the exterior Hilbert space and the
full adjoint in IGO27--28, including its original phase. For $r=1$
and $r=0$ the maps are the identities on the stated original spaces.
The original spectral sum $T=rc+iY$ intertwines on the exact domain
$\{F:\int |YF(Y)|^2d\nu(Y)<\infty\}$, as proved at the end of IGO.
These maps specify the original relation, the sum, and their retained
kernel throughout the receiving calculation.

\subsection{Return to the original arithmetic criterion}